\chapter{Pensiones}

\marginal{4.53\% del PIB}
Las pensiones y jubilaciones de la clasificación económica suman
\textbf{1,637,665.1 millones de pesos} en 2025, contra 1,499,038.6 en 2024: un
aumento real de \textbf{4.7\%} en un presupuesto cuyo gasto programable cae 4.1\%.
Pasan de 4.36 a \textbf{4.53\% del PIB} y equivalen al \textbf{17.6\% del gasto
neto total}.

\section{El perímetro, adjudicado y no elegido}

Este agregado no se puede tomar de un cuadro: hay que construirlo, y la
construcción admite error por doble contabilidad. El Gobierno Federal transfiere
a los institutos el dinero con el que pagan pensiones, y los institutos lo pagan;
sumar las dos cosas duplica 1,148,026.4 millones de pesos.

El perímetro que este documento usa es el tipo de gasto 4 de las entidades de
control directo más el tipo de gasto 4 del Gobierno Federal excluyendo la partida
45203, que es justamente esa transferencia. \textbf{La construcción no se defiende
por argumento sino por adjudicación}: el Anexo 3 del decreto publica los gastos
obligatorios con y sin pensiones, y su diferencia da 1,637,665.1 en 2025 y
1,499,038.6 en 2024, exactamente las dos cifras de arriba. Una fuente
independiente de los analíticos confirma el perímetro al millón de pesos en los
dos años.

\begin{table}[htbp]
\centering
\caption{Pensiones y jubilaciones por institución (mdp)}
\label{tab:C5_1}
\small
\begin{tabular}{lSSSSS}
\toprule
{Concepto} & {2024} & {2025} & {Dif. nominal} & {Dif. real} & {Var. real \%} \\
\midrule
{Entidades de control directo (pago directo)} & \num{1343107.4} & \num{1505658.9} & \num{162551.4} & \num{104797.8} & \num{7.5} \\
{Gobierno Federal, pago directo (sin 45203)} & \num{155931.2} & \num{132006.3} & \num{-23924.9} & \num{-30630.0} & \num{-18.8} \\
{CONSOLIDADO (clasificacion economica)} & \num{1499038.6} & \num{1637665.1} & \num{138626.5} & \num{74167.8} & \num{4.7} \\
{Partida informativa: transferencias del GF a las entidades (45203)} & \num{1009414.4} & \num{1148026.4} & \num{138612.0} & \num{95207.2} & \num{9.0} \\
{... GYN Instituto de Seguridad y Servicios Sociales de los Trabajadores del Estado} & \num{342043.9} & \num{390767.1} & \num{48723.2} & \num{34015.3} & \num{9.5} \\
{... GYR Instituto Mexicano del Seguro Social} & \num{870981.7} & \num{966726.3} & \num{95744.5} & \num{58292.3} & \num{6.4} \\
{... TVV CFE Consolidado} & \num{55961.3} & \num{64337.9} & \num{8376.5} & \num{5970.2} & \num{10.2} \\
{... TYY Pemex Consolidado} & \num{74120.5} & \num{83827.6} & \num{9707.2} & \num{6520.0} & \num{8.4} \\
\bottomrule
\end{tabular}
\par\vspace{2pt}\footnotesize Fuente: elaboración propia del ITED con los analíticos del PEF aprobado 2024 y 2025. Las diferencias en millones de pesos se reportan dos veces, nominal y real; la real está en pesos de 2025 con el deflactor de 4.3\% que declara el CGPE.
\end{table}


\section{Quién paga qué}

Las entidades de control directo pagan 1,505,658.9 millones de pesos, 7.5\% más en
términos reales. El Instituto Mexicano del Seguro Social pasa de 870,981.7 a
\textbf{966,726.3} (+6.4\% real), el Instituto de Seguridad y Servicios Sociales
de los Trabajadores del Estado de 342,043.9 a \textbf{390,767.1} (+9.5\%), la
Comisión Federal de Electricidad de 55,961.3 a 64,337.9 (+10.2\%) y Petróleos
Mexicanos de 74,120.5 a 83,827.6 (+8.4\%).

El pago directo del Gobierno Federal, en cambio, \textbf{cae 18.8\% real}, de
155,931.2 a 132,006.3. La composición se desplaza hacia los institutos.

\section{Las pensiones no contributivas}

Fuera de la clasificación económica de pensiones viven los programas del Ramo 20:
la Pensión para el Bienestar de las Personas Adultas Mayores (483,427.6 millones
de pesos en 2025), la Pensión para Personas con Discapacidad (28,961.4) y
\textbf{la Pensión Mujeres Bienestar, que no existía en 2024 y vale 15,000.0}. En
conjunto pasan de 492,909.0 a 527,389.1, un aumento real de 2.6\%.

\textbf{Sumadas a las contributivas, el compromiso pensionario del paquete es de
2,165,054.1 millones de pesos, 5.99\% del PIB.} Es la partida más grande del
presupuesto por cualquier definición.

\section{La comparación estructural: las cuotas y lo que pagan}

La razón canónica de esta casa para pensiones es \emph{pensiones del IMSS sobre
cuotas del IMSS}, con numerador en el tipo de gasto 4 del instituto y denominador
en el renglón 2.22 de la Ley de Ingresos aprobada. Los cuatro puntos son aprobado
contra aprobado, sin línea híbrida:

\begin{itemize}
\item 2022: 636,461.8 sobre 411,852.5 = \textbf{1.545}
\item 2023: 750,252.1 sobre 470,845.4 = \textbf{1.594}
\item 2024: 870,981.7 sobre 535,254.7 = \textbf{1.627}
\item 2025: 966,726.3 sobre 603,077.9 = \textbf{1.603}
\end{itemize}

\marginal{la razón baja por primera vez}
\textbf{Por primera vez en la serie, la razón baja.} No porque las pensiones
crezcan menos, sino porque las cuotas crecen más: 12.7\% nominal contra 11.0\% de
las pensiones del instituto. Un año no hace tendencia y conviene decirlo así. Lo
que sí se puede afirmar es que la brecha se abrió durante tres ejercicios
consecutivos y en 2025 se cerró dos centésimas y media, y que ninguna reforma
explica el movimiento: es composición del empleo formal contra maduración del
régimen anterior.

La serie \emph{ex ante}, que compara la exposición de motivos del proyecto contra
la iniciativa de ingresos, sigue congelada en sus dos puntos de 2020 y 2021.
\textbf{No se puede continuar porque la exposición de motivos del proyecto de
presupuesto no existe desde el ejercicio 2022}: está caída en el servidor de la
Secretaría, no retirada. Es la única serie de este documento que no avanza.

\section{Horizonte}

Es el único rubro del documento con horizonte oficial. El CGPE proyecta pensiones
y jubilaciones de \textbf{4.5\% del PIB en 2025 a 4.8\% en 2030}, tres décimas en
seis años. La proyección es notablemente plana para una población que envejece, y
el paquete no publica el supuesto demográfico que la sostiene. El capítulo 13
vuelve sobre esto.

\clearpage
